\documentclass{beamer}

\usetheme{Berlin}
\usecolortheme{Orchid}
\useoutertheme{miniframes}
\setbeamertemplate{navigation symbols}{}

\usepackage[utf8]{inputenc}
\usepackage[T1]{fontenc}
\usepackage{amsmath}
\usepackage{amssymb}
\usepackage{booktabs}
\usepackage{graphicx}
\graphicspath{{../images/}}

\usepackage{xcolor}
\usepackage{listings}
\usepackage{hyperref}
\usepackage{tikz}
\usetikzlibrary{arrows.meta,positioning,shapes.geometric,calc}

\lstset{
  language=Java,
  basicstyle=\ttfamily\small,
  keywordstyle=\color{blue},
  commentstyle=\color{gray},
  stringstyle=\color{teal},
  showstringspaces=false,
  columns=fullflexible,
  breaklines=true
}

% Per-slide source line (references shown on the slide itself).
\newcommand{\slidesources}[1]{%
  \vfill
  {\tiny\color{gray}\raggedright\sloppy\parbox{\linewidth}{Sources: #1}\par}%
}

% Unit-wise source strings for this subject (used by slides/notes).
% Path is relative to the lecture `latex/` folder (the compilation CWD).
% OOPS with Java (CSEG1044) — unit-wise sources
%
% These are short, student-facing reference strings intended to be used with:
%   - \slidesources{...}  (in beamer slides)
%   - \notesources{...}   (in lecture notes)
%
% Style inspiration: other subjects in My-Subjects/ use semicolon-separated sources
% and include an "accessed" date for web documentation.

\newcommand{\OOPSUnitISources}{Schildt, \textit{Java: A Beginner's Guide} (10e), McGraw-Hill (Java fundamentals + OOP basics).; Oracle Java Tutorials: Getting Started + Language Basics + Classes and Objects (accessed \today).; Java SE API docs (e.g., \texttt{String}, \texttt{Scanner}) (accessed \today).}

\newcommand{\OOPSUnitIISources}{Schildt, \textit{Java: The Complete Reference} (12e), McGraw-Hill (classes, inheritance, packages, interfaces).; Bloch, \textit{Effective Java} (3e), Addison-Wesley, 2018 (design choices; interfaces vs abstract classes).; Oracle Java Tutorials: Classes and Objects + Interfaces + Packages (accessed \today).; Java SE API docs (e.g., \texttt{Comparable}, \texttt{Runnable}) (accessed \today).}

\newcommand{\OOPSUnitIIISources}{Schildt, \textit{Java: The Complete Reference} (12e) (nested/inner classes, exceptions, threads, I/O).; Oracle Java Tutorials: Nested Classes + Exceptions + Concurrency + Basic I/O (accessed \today).; Java SE API docs (e.g., \texttt{Thread}, \texttt{AutoCloseable}) (accessed \today).}

\newcommand{\OOPSUnitIVSources}{Schildt, \textit{Java: The Complete Reference} (12e) (generics, lambdas, Swing, JDBC).; Bloch, \textit{Effective Java} (3e) (generics + lambdas).; Oracle Java Tutorials: Generics + Lambda Expressions + Swing + JDBC Basics (accessed \today).; Java SE API docs (e.g., \texttt{java.util.function}, \texttt{javax.swing}, \texttt{java.sql}) (accessed \today).}

\newcommand{\OOPSUnitVSources}{Schildt, \textit{Java: The Complete Reference} (12e) (Collections Framework + wrapper classes).; Oracle Java docs: Collections Framework overview (accessed \today).; Java SE API docs (\texttt{java.util} collections; wrapper classes in \texttt{java.lang}) (accessed \today).}

\newcommand{\OOPSUnitVISources}{Git SCM: \textit{Pro Git} (2e) (version control workflow), accessed \today.; GitHub Docs (repositories, issues, pull requests), accessed \today.; Oracle Java Tutorials: Swing + JDBC (accessed \today).; JUnit 5 User Guide (accessed \today).; Apache Maven Project: Getting Started (accessed \today).; Gradle User Manual: Getting Started (accessed \today).; UML class diagrams: UML 2.x overview/spec (accessed \today).}

\newcommand{\OOPSMiscWhyInterfacesSources}{Schildt, \textit{Java: The Complete Reference} (12e) (interfaces).; Bloch, \textit{Effective Java} (3e) (prefer interfaces where appropriate).; Oracle Java Tutorials: Interfaces (accessed \today).; Java SE API docs (e.g., \texttt{Comparable}, \texttt{Runnable}) (accessed \today).}



\title[OOPS with Java]{Object Oriented Programming (CSEG1044)}
\subtitle{Unit 02 -- Lecture 01: Classes, Objects, and Encapsulation (Basics)}
\author{Tofik Ali}
\institute{School of Computer Science, UPES Dehradun}
\date{\today}

\begin{document}

\begin{frame}[fragile]
  \titlepage
  \vspace{0.5em}
  \slidesources{\OOPSUnitIISources}
\end{frame}

\begin{frame}{Quick Links}
  \centering
  \hyperlink{sec:overview}{\beamerbutton{Overview}}\hspace{0.6em}
  \hyperlink{sec:concepts}{\beamerbutton{Key Topics}}\hspace{0.6em}
  \hyperlink{sec:demo}{\beamerbutton{Demo}}\hspace{0.6em}
  \hyperlink{sec:summary}{\beamerbutton{Summary}}
  \slidesources{\OOPSUnitIISources}
\end{frame}

\begin{frame}{Agenda}
  \tableofcontents
  \slidesources{\OOPSUnitIISources}
\end{frame}

\section{Overview}
\label{sec:overview}

\begin{frame}{Learning Outcomes}
  \begin{itemize}[<+->]
    \item Identify fields vs methods and explain state vs behavior.
    \item Explain references and aliasing using a small example.
    \item Apply encapsulation using \texttt{private} fields and validation methods.
  \end{itemize}
  \slidesources{\OOPSUnitIISources}
\end{frame}

\section{Key Topics}
\label{sec:concepts}

\begin{frame}{Unit 02: Classes, Inheritance, Packages and Interfaces}
  \begin{itemize}[<+->]
    \item Lecture focus: Classes, Objects, and Encapsulation (Basics)
  \end{itemize}
  \slidesources{\OOPSUnitIISources}
\end{frame}

\begin{frame}{Today: Roadmap}
  \begin{itemize}[<+->]
    \item Class fundamentals: fields, methods, state vs behavior
    \item Objects and reference semantics; object diagrams (mental model)
    \item Encapsulation motivation (invariants) + intro to access modifiers
  \end{itemize}
  \slidesources{\OOPSUnitIISources}
\end{frame}

\begin{frame}{Class vs Object (quick)}
  \begin{itemize}[<+->]
    \item \textbf{Class} = blueprint (type): fields + methods
    \item \textbf{Object} = instance at runtime (identity + state + behavior)
    \item State = values in fields; Behavior = what methods do
  \end{itemize}
  \slidesources{\OOPSUnitIISources}
\end{frame}

\begin{frame}{References and aliasing (important)}
  \begin{itemize}[<+->]
    \item In Java, variables of class type hold \textbf{references}.
    \item Assignment copies the reference (not the whole object).
    \item If two references point to the same object, changes are shared (\textbf{aliasing}).
  \end{itemize}
  \begin{block}{Why it matters}
    Many bugs come from unexpected shared updates.
  \end{block}
  \slidesources{\OOPSUnitIISources}
\end{frame}

\begin{frame}{Encapsulation = protect invariants}
  \begin{itemize}[<+->]
    \item Make fields \texttt{private}.
    \item Update state via methods that validate inputs.
    \item Keep object invariants always true (e.g., \texttt{0 <= cgpa <= 10}).
  \end{itemize}
  \begin{block}{Common mistake}
    Adding setters for everything and allowing invalid values.
  \end{block}
  \slidesources{\OOPSUnitIISources}
\end{frame}


\section{Demo / Example}
\label{sec:demo}

\begin{frame}[fragile]{Example: \texttt{Student} (encapsulation) (1/2)}
\begin{lstlisting}
public class Student {
    private final String rollNo;
    private final String name;
    private double cgpa; // invariant: 0..10

    public Student(String rollNo, String name, double cgpa) {
        if (cgpa < 0.0 || cgpa > 10.0)
            throw new IllegalArgumentException("cgpa");
        this.rollNo = rollNo;
        this.name = name;
        this.cgpa = cgpa;
    }
\end{lstlisting}
  \slidesources{\OOPSUnitIISources}
\end{frame}

\begin{frame}[fragile]{Example: \texttt{Student} (encapsulation) (2/2)}
\begin{lstlisting}
// continued...

    public void updateCgpa(double newCgpa) {
        if (newCgpa < 0.0 || newCgpa > 10.0)
            throw new IllegalArgumentException("cgpa");
        cgpa = newCgpa;
    }

    public double getCgpa() { return cgpa; }
}
\end{lstlisting}
  \slidesources{\OOPSUnitIISources}
\end{frame}

\begin{frame}[fragile]{Aliasing demo (observe)}
\begin{lstlisting}
public class Demo {
    public static void main(String[] args) {
        Student a = new Student("21CS1001", "Asha", 8.2);
        Student b = a;   // same object
        b.updateCgpa(9.0);
        System.out.println(a.getCgpa()); // prints 9.0
    }
}
\end{lstlisting}
  \slidesources{\OOPSUnitIISources}
\end{frame}

\begin{frame}{Mini Exercises (2--3 mins)}
  \begin{enumerate}[<+->]
    \item What is the difference between a class and an object?
    \item If \texttt{b = a} and \texttt{b} modifies the object, does \texttt{a} see it?
    \item Give one invariant for \texttt{BankAccount} or \texttt{LibraryBook}.
  \end{enumerate}
  \slidesources{\OOPSUnitIISources}
\end{frame}


\section{Summary}
\label{sec:summary}

\begin{frame}{Key Takeaways}
  \begin{itemize}[<+->]
    \item Classes define state + behavior; objects are runtime instances.
    \item Java object variables are references; aliasing is common.
    \item Encapsulation protects invariants using \texttt{private} + validation methods.
  \end{itemize}
  \slidesources{\OOPSUnitIISources}
\end{frame}

\begin{frame}{Checkpoint Questions}
  \begin{enumerate}[<+->]
    \item Why do we keep fields \texttt{private}?
    \item What is aliasing and why can it be risky?
    \item What is an invariant? Give one example.
  \end{enumerate}
  \slidesources{\OOPSUnitIISources}
\end{frame}

\begin{frame}{Next Lecture Preview}
  \textbf{Next:} Constructors, \texttt{this}, Overloading, \texttt{static}, and GC overview
  \vspace{0.6em}
  \begin{itemize}[<+->]
    \item We will learn how objects are created and how constructors enforce invariants.
    \item We will also learn when to use \texttt{static}.
  \end{itemize}
  \slidesources{\OOPSUnitIISources}
\end{frame}

\end{document}
